%%==========================================================================================
%%
%%   POSTER TEMPLATE  +  USAGE GUIDE  for  beamerthemePoster.sty
%%
%%   This file is two things at once:
%%     1. A working main.tex you can compile right now (it has no external image/figure/
%%        bib dependencies, so `pdflatex` will produce a poster on the first try).
%%     2. A commented walkthrough of every feature beamerthemePoster.sty exposes, so you
%%        can use it as a reference the next time you start a new poster.
%%
%%   DIRECTORY LAYOUT THIS FILE ASSUMES
%%   -----------------------------------------------------------------------------------
%%     project/
%%       poster/
%%         main.tex                      <- this file
%%         beamerthemePoster.sty         <- the layout engine (headline/block/footline)
%%         beamercolorthemeUofI.sty      <- the color palette (swap for your own brand)
%%         references.bib                <- optional, only needed if you cite something
%%       _mathematics.tex                <- optional personal math-macro library
%%       _definitions.tex                <- optional personal math-macro library
%%       figures/                        <- recommended home for reusable TikZ figures
%%                                          (shared between paper and poster), each as
%%                                          its own .tex file you \input{} from a block
%%
%%   QUICK START
%%   -----------------------------------------------------------------------------------
%%     1. Edit the title/author/institute block below.
%%     2. Replace the example block content in each column with your own.
%%     3. (Optional) Re-enable \logo, \qrcode, \footercontent, and the bibliography once
%%        you have the corresponding image/bib files -- see the comments at each command.
%%     4. Compile with:  pdflatex main.tex  (run biber + pdflatex again if you use \cite)
%%
%%   WHAT'S ALREADY LOADED FOR YOU
%%   -----------------------------------------------------------------------------------
%%     beamerthemePoster.sty already \RequirePackage{}'s a large standard toolkit, so you
%%     do NOT need to \usepackage{} any of the following yourself (doing so will trigger
%%     "already defined" warnings): amsmath/amsfonts/amssymb/mathtools, tikz, pgfplots,
%%     booktabs, arydshln, multirow/multicol, tcolorbox, algorithm/algpseudocode, soul,
%%     siunitx, csquotes, caption/subcaption, float, nameref, dsfont, upgreek, nicefrac,
%%     cancel, systeme, biblatex (backend=biber). If you need something not in this list,
%%     load it here in main.tex, after \usetheme.
%%
%%==========================================================================================


%%------------------------------------------------------------------------------------------
%%   1. DOCUMENT CLASS & PHYSICAL POSTER SIZE
%%------------------------------------------------------------------------------------------
%%   beamerposter's `size=custom` lets you set the exact printed dimensions in cm.
%%   Convert inches to cm with: cm = inches * 2.54.
%%     48 in x 36 in  ->  121.92 cm x 91.44 cm   (landscape, used below)
%%     36 in x 48 in  ->   91.44 cm x 121.92 cm  (portrait: just swap width/height)
%%     A0             ->  118.9 cm  x  84.1 cm
%%   `scale` blows up every font size and line width uniformly without changing the
%%   physical page size -- it's how you make text/figures read well from a few feet away.
%%   1.4-1.8 is a good starting range; bump it up if a draft print looks too small.
\documentclass[final,dvipsnames]{beamer}
\usepackage[size=custom, width=121.92, height=91.44, scale=1.5]{beamerposter}


%%------------------------------------------------------------------------------------------
%%   2. PERSONAL MACRO LIBRARIES (optional)
%%------------------------------------------------------------------------------------------
%%   These two files are reusable math-notation shortcuts (e.g. \sfA -> \mathsf{A}) that
%%   are independent of the poster theme -- keep them, delete them, or point the \input
%%   path at your own macro files. If you don't have any, just delete these two lines and
%%   write your math macros directly below instead.
% \newcommand{\composition}{\circ}
\newcommand{\define}{=\vcentcolon}
\newcommand{\defined}{\vcentcolon=}
\newcommand{\isomorphic}{\cong}
\newcommand{\complexes}{\mathbb{C}}
\newcommand{\reals}{\mathbb{R}}
\newcommand{\integers}{\mathbb{Z}}
\newcommand{\naturals}{\mathbb{N}}
\newcommand{\continuous}{C}
\newcommand{\euclideans}{\mathbb{E}}
\newcommand{\Mat}[1][]{\ifthenelse{\equal{#1}{}}{\textnormal{Mat}}{\textnormal{Mat}(#1)}}
\newcommand{\SO}[1]{\textnormal{SO}({#1})}
\newcommand{\SE}[1]{\textnormal{SE}({#1})}
% \newcommand{\so}[1]{\mathfrak{so}(#1)}
\newcommand{\se}[1]{\mathfrak{se}(#1)}
\newcommand{\RSO}[1]{R_{\textnormal{\tiny{$\SO$}}}\hspace{-3pt}\left(#1\right)}
\newcommand{\RSE}[1]{R_{\textnormal{\tiny{$\SE$}}}\hspace{-3pt}\left(#1\right)}
\newcommand{\tangent}[1]{
    \ifthenelse{\equal{#1}{}}
    {{T}}
    {{T_{#1}}}
}
\newcommand{\dualtangent}[1]{
    \ifthenelse{\equal{#1}{}}
    {{T^*}}
    {{T_{#1}^*}}
}

\newcommand{\leftaction}[1]{L_{#1}}
\newcommand{\leftextension}[2][e]{\tangent{#1}\leftaction{#2}}
\newcommand{\dualleftextension}[2][e]{\dualtangent{#1}\leftaction{#2}}
\newcommand{\leftextensionfromvector}[1]{\Xi_{#1}}
\newcommand{\MaurerCartanform}[1]{\chi_{#1}}
% \newcommand{\TSE}[2]{
%     \ifthenelse{\equal{#1}{}}
%     {{T\SE{#2}}}
%     {{T_{#1}\SE{#2}}}
% }
% \newcommand{\TSE}[2]{
%     \tangent[#1]\SE{#2}
% }
% \newcommand{\dualTSE}[2][]{
%     \ifthenelse{\equal{#1}{}}
%     {{\dualtangent\SE{#2}}}
%     {{\dualtangent[#1]\SE{#2}}}
% }
\newcommand{\dualse}[1]{\mathfrak{se}^*(#1)}
\newcommand{\sebasis}{X}
\newcommand{\Adjoint}[1]{
    \ifthenelse{\equal{#1}{}}
    {\textnormal{Ad}}
    {\textnormal{Ad}_{#1}}
}

\newcommand{\adjoint}[1]{
    \ifthenelse{\equal{#1}{}}
    {\textnormal{ad}}
    {\textnormal{ad}_{#1}}
}
\newcommand{\coAdjoint}[1]{
    \ifthenelse{\equal{#1}{}}
    {\textnormal{Ad}^*}
    {\textnormal{Ad}^*_{#1}}
}

\newcommand{\coadjoint}[1]{
    \ifthenelse{\equal{#1}{}}
    {\textnormal{ad}^*}
    {\textnormal{ad}^*_{#1}}
}
\newcommand{\maptoliealgebra}[1]{\left(#1\right)_{\wedge}}
\newcommand{\maptovector}[1]{\left(#1\right)_{\vee}}
\newcommand{\maptodualliealgebra}[1]{\left(#1\right)^*_{\wedge}}

% \newcommand{\maptoliealgebra}[1]{\left(#1\right)_{\times}}
% \newcommand{\maptovector}[1]{\left(#1\right)_{\vert}}

\newcommand{\perpendicular}{\bot}

\newcommand{\circconvolute}[1][\ast]{\raisebox{1pt}{\scriptsize{\textcircled{\tiny{$#1$}}}}}
\newcommand{\convolute}{*}

\newcommand{\dif}{\mathrm{d}}
\newcommand{\Dif}{\mathrm{D}}
\newcommand{\gradient}{\nabla}
\newcommand{\divergence}{\nabla\cdot}
\newcommand{\Hessian}{\mathbf{H}}

\newcommand{\zeros}{\mathbf{0}}
\newcommand{\ones}{\mathbf{1}}
\newcommand{\Set}[1]{\left\{#1\right\}}
\newcommand{\inlineset}[1]{\{#1\}}
\newcommand{\given}{\,\middle|\,}
\newcommand{\at}{|}
% \renewcommand{\exp}[1]{\textnormal{ exp}\left(#1\right)}
\newcommand{\csch}[1]{\textnormal{ csch}\left(#1\right)}
\newcommand{\sech}[1]{\textnormal{ sech}\left(#1\right)}
\newcommand{\softmax}[1]{\textnormal{softmax}\left(#1\right)}
\newcommand{\inlinesoftmax}[1]{\textnormal{softmax}(#1)}

\newcommand{\inline}[1]{%
  \begingroup
    \let\left\relax
    \let\right\relax
    #1%
  \endgroup
}
\newcommand{\sgn}[1]{\textnormal{sgn}\left(#1\right)}
\newcommand{\abs}[1]{\left|#1\right|}
\newcommand{\norm}[1]{\left\lVert#1\right\rVert}
\renewcommand{\det}[1]{\textnormal{det}\left(#1\right)}
\newcommand{\inverse}{{-1}}
\newcommand{\pseudoinverse}{+}
\newcommand{\conjugate}{*}
\newcommand{\transpose}{\intercal}
\newcommand{\conjugatetranspose}{\dagger}
\newcommand{\Trace}[1]{\textnormal{Tr}\left(#1\right)}
\newcommand{\diag}[1]{\textnormal{diag}\left\{#1\right\}}
\newcommand{\inner}[2]{\left\langle #1, #2 \right\rangle}
\newcommand{\pairing}[2]{\left\llangle #1, #2 \right\rrangle}
\newcommand{\liebracket}[2]{\left[ #1, \ #2 \right]}
\newcommand{\Frobenius}{\textnormal{F}}
\newcommand{\Frobeniusnorm}[1]{\norm{#1}_\Frobenius}
\newcommand{\distance}[2]{\textnormal{dist}(#1, #2)}
% \newcommand{\Vector}[1]{\boldsymbol{#1}}
% \newcommand{\Matrix}[1]{
% 			\ifthenelse{\equal{#1}{\kappa}}{\boldsymbol{\kappa}}{
% 			\ifthenelse{\equal{#1}{\nu}}{\boldsymbol{\nu}}{
% 			\ifthenelse{\equal{#1}{\omega}}{\boldsymbol{\omega}}{
% 			\ifthenelse{\equal{#1}{\sigma}}{\boldsymbol{\sigma}}{
% 			\ifthenelse{\equal{#1}{\theta}}{\boldsymbol{\theta}}{
% 			\ifthenelse{\equal{#1}{\xi}}{\boldsymbol{\xi}}{
%			\ifthenelse{\equal{#1}{\lambda}}{\boldsymbol{\lambda}}{
%			\mathbf{#1}}}}}}}}
% 			}
			
\newcommand{\Vector}[1]{{#1}}
\newcommand{\Matrix}[1]{\mathrm{#1}}
\newcommand{\identity}[1]{
    \ifthenelse{\equal{#1}{}}
    {\Matrix{I}}
    {\Matrix{I}_{#1}}
}

\newcommand{\union}{\bigcup}

\newcommand{\ProbabilitySimplex}{\Delta}
\newcommand{\ProbabilitySpace}[1]{\mathcal{P}\left(#1\right)}
\newcommand{\sigmaAlgebra}[1]{\sigma\left(#1\right)}
\newcommand{\indicator}[1]{\textbf{1}_{\left\{#1\right\}}}
\NewDocumentCommand{\Probability}{o m}{
  \mathsf{P}
  \IfValueT{#1}{_{#1}}
  \IfNoValueT{#1}{}
  \left(#2\right)
}
\NewDocumentCommand{\Expectation}{o m}{
  \mathsf{E}
  \IfValueT{#1}{_{#1}}
  \IfNoValueT{#1}{}
  \left(#2\right)
}
\NewDocumentCommand{\Variance}{o m}{
  \mathsf{Var}
  \IfValueT{#1}{_{#1}}
  \IfNoValueT{#1}{}
  \left(#2\right)
}
\DeclareMathOperator*{\logsumexp}{logsumexp}
\newcommand{\entropy}[1]{\mathsf{H}\left(#1\right)}

\newcommand{\Dim}{\textnormal{Dim}}
\renewcommand{\dim}{\mathsf{d}}

\DeclareMathOperator*{\argmax}{arg\,max}
\DeclareMathOperator*{\argmin}{arg\,min}

\newcommand{\subjectto}{\textnormal{s.t.}}

\newcommand{\Order}[1]{O\left(#1\right)}

\newcommand{\mysetminusD}{\hbox{\tikz{\draw[line width=0.6pt,line cap=round] (3pt,0) -- (0,6pt);}}}
\newcommand{\mysetminusT}{\mysetminusD}
\newcommand{\mysetminusS}{\hbox{\tikz{\draw[line width=0.45pt,line cap=round] (2pt,0) -- (0,4pt);}}}
\newcommand{\mysetminusSS}{\hbox{\tikz{\draw[line width=0.4pt,line cap=round] (1.5pt,0) -- (0,3pt);}}}

\renewcommand{\setminus}{\mathbin{\mathchoice{\mysetminusD}{\mysetminusT}{\mysetminusS}{\mysetminusSS}}}
% \newcommand{\raisedtri}[1]{%
  % 0.3ex is how high to raise it. 0.7 is the size (70%).
  % #1 passes the math style (display, text, etc.) so it shrinks in subscripts
  \raisebox{0.2ex}{\scalebox{0.6}{$#1\vartriangle$}}%
}

\makeatletter
\@tfor\letter:=abcdefghijklmnopqrstuvwxyz\do{%
  % \tt\letter
  \expandafter\edef\csname tt\letter\endcsname{\noexpand\textnormal{\letter}}%
  % \sf\letter
  \expandafter\edef\csname sf\letter\endcsname{\noexpand\mathsf{\letter}}%
  % \rm\letter
  \expandafter\edef\csname rm\letter\endcsname{\noexpand\mathrm{\letter}}%
  % \delta\letter
  \expandafter\edef\csname delta\letter\endcsname{%
    \noexpand\mathchoice
      {\noexpand\raisedtri\displaystyle\noexpand\mkern-1mu \letter}%
      {\noexpand\raisedtri\textstyle\noexpand\mkern-1mu \letter}%
      {\noexpand\raisedtri\scriptstyle\noexpand\mkern-1mu \letter}%
      {\noexpand\raisedtri\scriptscriptstyle\noexpand\mkern-1mu \letter}%
  }%
}
\@tfor\letter:=ABCDEFGHIJKLMNOPQRSTUVWXYZ\do{%
  % \tt\letter
  \expandafter\edef\csname tt\letter\endcsname{\noexpand\textnormal{\letter}}%
  % \cl\letter
  \expandafter\edef\csname cl\letter\endcsname{\noexpand\mathcal{\letter}}%
  % \sf\letter
  \expandafter\edef\csname sf\letter\endcsname{\noexpand\mathsf{\letter}}%
  % \bb\letter
  \expandafter\edef\csname bb\letter\endcsname{\noexpand\mathbb{\letter}}%
  % \rm\letter
  \expandafter\edef\csname rm\letter\endcsname{\noexpand\mathrm{\letter}}%
  % \delta\letter
  \expandafter\edef\csname delta\letter\endcsname{%
    \noexpand\mathchoice
      {\noexpand\raisedtri\displaystyle\noexpand\mkern-1mu \letter}%
      {\noexpand\raisedtri\textstyle\noexpand\mkern-1mu \letter}%
      {\noexpand\raisedtri\scriptstyle\noexpand\mkern-1mu \letter}%
      {\noexpand\raisedtri\scriptscriptstyle\noexpand\mkern-1mu \letter}%
  }%
}
\makeatother

% negative kern that pulls a subscript back toward its base after \overline or \bar,
% which box the base and drop the italic correction, e.g. \bar{Y}_{\sss t}
\newcommand{\sss}{\mkern-4.5mu} % slanted subscript spacing


%%------------------------------------------------------------------------------------------
%%   3. THEME + COLOR THEME
%%------------------------------------------------------------------------------------------
%%   \usetheme[...]{Poster} accepts four length-valued options, all optional. Defaults
%%   shown below; override only the ones you need to change.
%%
%%     headerHeight     (default 10cm)        Height of the top headline bar. Also used
%%                                             to auto-scale \logo and \qrcode (see §6):
%%                                             whatever image you pass is resized so its
%%                                             height equals headerHeight, and the title
%%                                             column shrinks to make room automatically.
%%     paddingWidth     (default 2em)         Vertical breathing room added above/below
%%                                             the headline and footline content, and part
%%                                             of the horizontal margin on those two bars.
%%     blockTitleHeight (default 6\baselineskip)
%%                                             Fixed height of every block's title bar, so
%%                                             titles of different lengths still line up
%%                                             across columns. Increase if your longest
%%                                             block title wraps to two lines and looks cramped.
%%     sidePadding      (default 3.2em)       Extra horizontal margin between the headline/
%%                                             footline content and the poster's left/right
%%                                             edge (added on top of paddingWidth).
%%
%%   Note these are plain TeX lengths, so any unit works (cm, in, pt, em, \baselineskip...).
\usetheme[headerHeight=9cm]{Poster}

%%   The color theme only has to set a handful of beamer color names; copy
%%   beamercolorthemeUofI.sty to e.g. beamercolorthemeMyConf.sty and edit the RGB triples
%%   to rebrand for a different university/conference. The names the layout engine reads:
%%     headline                 (fg)      title/author/institute text color
%%     block title               (fg,bg)  block title bar text + background
%%     block body                 (fg)    body text color
%%     itemize item/subitem/subsubitem (fg)  bullet + bullet-point color
%%     footline                  (fg,bg)  footer bar text + background
\usecolortheme{UofI}


%%------------------------------------------------------------------------------------------
%%   4. EXTRA TIKZ/SYMBOL HELPERS (project-specific, so they live here, not in the theme)
%%------------------------------------------------------------------------------------------
%%   A named TikZ fading you can apply to a shape/shading so it fades out toward one edge,
%%   e.g. `\shade[path fading=fadeing, left color=uiucBlue] (0,0) rectangle (4,1);`
%%   Delete this if none of your figures need a fade effect.
\tikzfading[name=fadeing,
  left color=transparent!40,
  right color=transparent!0]

%%   pifont (loaded by the theme) gives you dingbat symbols by number; \ding{55} is a
%%   cross mark. Add a matching checkmark the same way if you need one: \ding{51}.
\newcommand{\xmark}{\ding{55}}%
\newcommand{\cmark}{\ding{51}}%


%%------------------------------------------------------------------------------------------
%%   5. BIBLIOGRAPHY (optional -- only needed if you use \cite/\autocite)
%%------------------------------------------------------------------------------------------
%%   biblatex (backend=biber, style=authoryear) is already configured by the theme; you
%%   only need to point it at your .bib file. The \IfFileExists guard below means this
%%   template still compiles cleanly even before you've created references.bib -- once
%%   you add the file, citations and \printbibliography (see the References block near
%%   the end of this file) start working with no other changes needed.
%%   Compile sequence once you have citations: pdflatex -> biber -> pdflatex -> pdflatex.
\IfFileExists{references.bib}{\addbibresource{references.bib}}{}


%%------------------------------------------------------------------------------------------
%%   6. TITLE, AUTHORS, INSTITUTE, LOGO, QR CODE, FOOTER
%%------------------------------------------------------------------------------------------
%%   Any of \author, \institute, \logo, \qrcode, \footercontent may simply be omitted --
%%   the headline/footline templates detect the missing piece and collapse its space
%%   instead of leaving a gap (e.g. no \logo call means the title block just centers in
%%   the full width with no logo column reserved).
\title{Your Poster Title Goes Here}
\author{First Author, Second Author, Third Author}
\institute{Your Department, Your University}

%%   \logo and \qrcode both auto-resize whatever image you give them to exactly
%%   headerHeight tall, then reserve precisely that much horizontal space for them --
%%   you never need to manually size these images. Any \includegraphics-compatible
%%   format works (png/jpg/pdf); vector PDF is sharpest at poster scale.
%%   Uncomment once you have the image files in this directory:
% \logo{\includegraphics{logo.png}}
\qrcode{\includegraphics{qr-code.png}}

%%   The footline bar is only rendered if \footercontent is called at all; comment the
%%   line below back out to remove the footer entirely. \hfill works inside it for a
%%   left/right-justified two-part footer, as shown.
\footercontent{\footnotesize Your Conference Name \& Year \hfill Location, Country\phantom{~}}


%%------------------------------------------------------------------------------------------
%%   7. COLUMN LAYOUT
%%------------------------------------------------------------------------------------------
%%   The body is built with beamer's native \begin{columns}[t]...\end{columns}, not a
%%   theme-specific column command, so any number of columns works. Below is the common
%%   3-column layout: three columns at 0.31\paperwidth each use 0.93\paperwidth in total;
%%   beamer's columns environment auto-distributes the remaining 0.07\paperwidth as gaps
%%   between/around the columns, so you don't normally need to insert spacing manually.
%%     - For 2 columns: use \colWidth = ~0.47\paperwidth, two \begin{column} blocks.
%%     - For 4 columns: use \colWidth = ~0.23\paperwidth, four \begin{column} blocks.
\newlength{\colWidth}
\setlength{\colWidth}{0.31\paperwidth}

%%   \sepwidth is not consumed automatically (beamer already distributes the gap for you,
%%   see above) -- it's provided here only as a convenience length in case you want manual
%%   control, e.g. `\begin{column}{\colWidth}...\end{column}\hspace{\sepwidth}\begin{column}{\colWidth}...`
\newlength{\sepwidth}
\setlength{\sepwidth}{0.1\paperwidth}


\begin{document}
  \begin{frame}[t]

    %% This negative \vspace compensates for the 0.25\baselineskip the theme's "block
    %% begin" template automatically adds above every block's title bar (see
    %% beamerthemePoster.sty), so the very first block in each column doesn't sit lower
    %% than the headline intends. Tune the amount to taste, or delete it if you'd rather
    %% have a touch more breathing room under the headline.
    \vspace{-1.5\baselineskip}

    \begin{columns}[t]

      %%================================================================================
      %% COLUMN 1
      %%================================================================================
      \begin{column}{\colWidth}

        \begin{block}{Block Title --- e.g. Problem Statement}

          %% --- Example figure: a small self-contained TikZ diagram -----------------
          %% In real projects it's worth keeping each figure in its own file under a
          %% sibling ../figures/ directory and \input{}-ing it here, so the exact same
          %% TikZ source is shared between your paper and your poster:
          %%   \input{../figures/my-figure.tex}
          %% Scaling the picture (and shrinking nodes to match) is the usual way to make
          %% a figure drawn for a paper column also look right inside a poster block.
          \begin{figure}[h]
            \centering
            \begin{tikzpicture}[
                scale=0.85, every node/.append style={scale=0.7},
                node distance=1.4cm,
                box/.style={draw, rounded corners, minimum width=2.6cm, minimum height=1cm, align=center},
                arrow/.style={-{Latex[length=2.2mm]}, thick}
              ]
              \node[box] (a) {Input};
              \node[box, right=of a] (b) {Process};
              \node[box, right=of b] (c) {Output};
              \draw[arrow] (a) -- (b);
              \draw[arrow] (b) -- (c);
            \end{tikzpicture}
            \vspace{-0.25\baselineskip}
          \end{figure}

          %% --- The \highlight environment: a colored tcolorbox for a takeaway claim --
          %% Its colors (currently uiucOrange!20!white) are fixed inside the theme file --
          %% it takes no optional arguments. To get a second highlight color, either edit
          %% the `highlight` environment definition in beamerthemePoster.sty directly, or
          %% duplicate it there under a new name (e.g. `highlightAlt`) with different colors.
          \begin{highlight}
            State the single most important takeaway of this block here, in one or two
            sentences, so it's readable from across the room.
          \end{highlight}

        \end{block}

        %% Blocks already add their own 0.25\baselineskip of space above and below; this
        %% small negative \vspace between blocks is a common trick to tighten up a poster
        %% that's running long on vertical space. Remove it if you'd rather keep the
        %% default spacing.
        \vspace*{-0.25\baselineskip}

        \begin{block}{Block Title --- e.g. Related Work}

          %% --- \paragraph{...}: an inline bold lead-in, NOT a new block ---------------
          %% Redefined by the theme to drop section numbering and just bold the word(s)
          %% you give it, with a small skip above and a space after -- handy for breaking
          %% one block into a few short labeled sub-points without the overhead of \heading.
          \paragraph{First Topic} A short description of the first related approach goes here.
          \begin{itemize}
            \item Strength or property worth calling out
            \item Limitation that motivates your approach
          \end{itemize}

          \paragraph{Second Topic} A short description of a second related approach.
          \begin{itemize}
            \item Strength or property worth calling out
            \item Limitation that motivates your approach
          \end{itemize}

        \end{block}

        \vspace*{-0.25\baselineskip}

        \begin{block}{Acknowledgements}
          This work is supported in part by [funding source/grant number]. We thank
          [collaborator] for [contribution].
        \end{block}

      \end{column}


      %%================================================================================
      %% COLUMN 2
      %%================================================================================
      \begin{column}{\colWidth}

        \begin{block}{Block Title --- e.g. Method}

          %% --- \heading{...}: an unboxed sub-heading inside a block -------------------
          %% Use this (rather than starting a new \block) when two sub-topics belong
          %% under one title bar but still deserve a visual break between them.
          \heading{Background}
          A sentence or two of background the reader needs before the main contribution.

          \heading{Main Idea}

          %% --- The \namelist environment: a definition-list for key facts ------------
          %% The argument to \begin{namelist}{...} should be the longest label you use,
          %% so all the bullet text aligns -- here that's "Objective:".
          \begin{namelist}{Objective:}
            \item[Input:] what goes into your method
            \item[Output:] what comes out
            \item[Objective:] the quantity you optimize, e.g. $\displaystyle\min_{\theta}\ \mathcal{L}(\theta)$
          \end{namelist}

        \end{block}

        \vspace*{-0.25\baselineskip}

        \begin{block}{Block Title --- e.g. Main Contribution}

          \begin{highlight}
            \paragraph{Key Result} A one-sentence statement of your main theoretical or
            empirical result, written so it stands on its own.
          \end{highlight}

          A short paragraph elaborating on the result, with a display equation if useful:
          \begin{equation*}
            f_\theta : \mathcal{X} \to \mathcal{Y}, \qquad x \mapsto f_\theta(x)
          \end{equation*}

        \end{block}

      \end{column}


      %%================================================================================
      %% COLUMN 3
      %%================================================================================
      \begin{column}{\colWidth}

        \begin{block}{Results \& Experiments}

          %% --- Example table: booktabs + multirow, already styled by the theme -------
          %% beamerthemePoster.sty thickens \heavyrulewidth/\lightrulewidth/\cmidrulewidth
          %% so booktabs rules stay visible at poster scale -- you don't need to set
          %% those yourself, just use \toprule/\midrule/\bottomrule/\cmidrule as normal.
          \begin{table}[h]
            \centering
            \renewcommand{\arraystretch}{1.25}
            \begin{tabular}{lccc}
              \toprule
              Method & Metric A & Metric B & Metric C \\
              \cmidrule(lr){1-1} \cmidrule(lr){2-4}
              \textbf{Ours} & \multirow{2}{*}{$0.12$} & \multirow{2}{*}{$0.98$} & fast \\
              Baseline & & & slow~\xmark \\
              \bottomrule
            \end{tabular}
          \end{table}

          \vspace*{0.375\baselineskip}

          %% --- Example chart: a self-contained pgfplots axis -------------------------
          %% As with the figure in Column 1, in a real project you'd usually keep this
          %% in ../figures/results-chart.tex and \input{} it here.
          \begin{tikzpicture}
            \begin{axis}[
                width=0.95\colWidth, height=0.45\colWidth,
                xlabel={Iteration}, ylabel={Loss},
                xtick distance=1,
                legend pos=north east, grid=major,
              ]
              \addplot+[mark=none, thick] coordinates {(0,2.5) (1,1.8) (2,1.3) (3,1.0) (4,0.8) (5,0.65)};
              \addlegendentry{Method A}
              \addplot+[mark=none, thick, dashed] coordinates {(0,2.6) (1,2.1) (2,1.7) (3,1.45) (4,1.25) (5,1.1)};
              \addlegendentry{Method B}
            \end{axis}
          \end{tikzpicture}

        \end{block}

        \vspace*{-0.25\baselineskip}

        \begin{block}{References}
          %% Mirrors the \IfFileExists guard in §5: shows a friendly placeholder until
          %% references.bib exists, then prints the real bibliography automatically.
          \IfFileExists{references.bib}{\printbibliography}{
            \textit{Add a \texttt{references.bib} file next to this one, then cite with
            \texttt{\textbackslash autocite\{key\}} -- this block will populate automatically.}
          }
        \end{block}

      \end{column}

    \end{columns}
  \end{frame}
\end{document}